\section{Enterprise Data Governance, Privacy Threat Model, and Economic TCO}
\label{sec:enterprise_governance}

A core impediment to enterprise adoption of Retrieval-Augmented Generation is the dual challenge of regulatory non-compliance and prohibitive commercial cloud API expenditures. This section provides a formal security threat model, a data governance compliance analysis, and an empirical Total Cost of Ownership (TCO) break-even model.

\subsection{Security and Data Privacy Threat Model}
We formalize the enterprise data residency threat model under the standard honest-but-curious and malicious exfiltration paradigms:

\begin{figure*}[t]
\centering
\begin{minipage}{0.96\textwidth}
\begin{verbatim}
====================================================================================================
                        ENTERPRISE ZERO-DATA-EGRESS SECURITY BOUNDARY
====================================================================================================
           CORPORATE INTRANET PERIMETER (Zero Public Internet Egress)
  ┌────────────────────────────────────────────────────────────────────────────────────────┐
  │                                                                                        │
  │   [User Client] ──(HTTPS/mTLS)──► [AMORA Orchestration Gateway (FastAPI)]              │
  │                                                   │                                    │
  │                     ┌─────────────────────────────┼─────────────────────────────┐      │
  │                     ▼                             ▼                             ▼      │
  │           [Local FAISS Vector]           [Internal Relational]         [Compliance PDF]│
  │           (MiniLM-L6 Embeddings)         (PostgreSQL / SQLite)         (Local Handbooks│
  │                     │                             │                             │      │
  │                     └─────────────────────────────┼─────────────────────────────┘      │
  │                                                   ▼                                    │
  │                                      [Evidence Fusion Engine]                          │
  │                                                   │                                    │
  │                                                   ▼                                    │
  │                                       [Local Qwen 2.5 7B LLM]                          │
  │                                         (Ollama on RTX 3060)                           │
  │                                                   │                                    │
  │                                                   ▼                                    │
  │                                      [Grounded Internal Answer]                        │
  │                                                                                        │
  └────────────────────────────────────────────────────────────────────────────────────────┘
                   X  BLOCKED: Zero Outbound Public WAN Connections  X
                    (No OpenAI API, No Anthropic API, No External Cloud)
====================================================================================================
\end{verbatim}
\end{minipage}
\caption{The AMORA Zero-Data-Egress Security Architecture. All intermediate vector representations, SQL queries, database credentials, and generative token decoding operate strictly within the private enterprise network boundary, eliminating data leakage vulnerabilities.}
\label{fig:security_boundary}
\end{figure*}

\subsubsection{Attacker Capabilities and Assets at Risk}
In enterprise knowledge operations, the primary assets requiring cryptographic and procedural isolation comprise:
\begin{enumerate}
    \item \textbf{Operational Relational Schemas ($\mathcal{M}_{\text{schema}}$):} Column identifiers, table relationships, and foreign key constraints that reveal corporate internal architecture.
    \item \textbf{Confidential Transactional Records ($\mathcal{D}_{\text{records}}$):} Financial ledgers, supplier price agreements, employee compensation details, and customer PII (Personally Identifiable Information).
    \item \textbf{Proprietary Intellectual Property ($\mathcal{D}_{\text{IP}}$):} Internal technical manuals, standard operating procedures, and unreleased product roadmaps.
    \item \textbf{User Prompt Queries ($q$):} High-intent executive inquiries that leak market strategy or pending legal disputes.
\end{enumerate}

Under commercial cloud-hosted RAG architectures (e.g., querying GPT-4o or Gemini APIs via public internet endpoints), every retrieved document chunk and user query must be transmitted across external wide-area network (WAN) boundaries. Even with Transport Layer Security (TLS 1.3), this architecture exposes the enterprise to three critical vulnerabilities:
\begin{itemize}
    \item \textbf{Third-Party Subprocessor Risk:} Cloud API providers routinely log prompt payloads for telemetry, model refinement, or compliance monitoring, violating strict corporate confidentiality covenants.
    \item \textbf{Man-in-the-Middle (MitM) and DNS Hijacking:} Transit across public routing nodes exposes high-value corporate communications to interception or state-level eavesdropping.
    \item \textbf{Regulatory Non-Compliance Liabilities:} Cross-border data transit violates sovereign data residency requirements.
\end{itemize}

\subsubsection{The AMORA Zero-Egress Architectural Invariant}
AMORA completely eliminates external data egress by adhering to a strict local containment invariant:
\begin{equation}
    \forall q \in \mathcal{Q}, \quad \text{PacketEgress}(q) \cap \text{PublicWAN} = \emptyset
\end{equation}
As illustrated in Figure \ref{fig:security_boundary}:
\begin{itemize}
    \item Dense vector embeddings are generated on-premises via local CPU execution of \texttt{all-MiniLM-L6-v2}.
    \item Vector nearest-neighbor indexing runs in local host memory via FAISS.
    \item SQL queries execute against local SQLite/PostgreSQL instances behind the internal DMZ.
    \item Autoregressive sequence generation executes on a dedicated, on-premises NVIDIA RTX 3060 GPU through Ollama.
    \item \textbf{Verifiable Auditability:} Network socket inspection confirms that during end-to-end benchmark execution across all 20 enterprise queries, the AMORA host emitted **zero outbound TCP/UDP packets** outside the local intranet subnet.

\end{itemize}

\subsection{Regulatory Compliance Mapping (GDPR, HIPAA, SOC 2)}
Deploying AMORA directly satisfies the stringent legal requirements imposed by leading global regulatory frameworks:

\subsubsection{EU General Data Protection Regulation (GDPR)}
\begin{itemize}
    \item \textbf{Article 22 (Automated Decision-Making and Profiling):} AMORA's Evidence Fusion engine produces auditable, deterministic provenance tags for every retrieved fact, enabling complete human-in-the-loop auditability of automated outputs.
    \item \textbf{Article 28 (Processor Obligations) and Article 32 (Security of Processing):} By eliminating third-party cloud API subprocessors, enterprises avoid the legal necessity of third-party Data Processing Agreements (DPAs) and eliminate risks associated with subprocessor data breaches.
    \item \textbf{Article 44--49 (Cross-Border Data Transfers):} Because all processing occurs within the enterprise's local jurisdiction, no international data transfers occur, eliminating invalidation risks under Schrems II rulings.
\end{itemize}

\subsubsection{Health Insurance Portability and Accountability Act (HIPAA)}
\begin{itemize}
    \item \textbf{Security Rule (§ 164.312):} AMORA preserves Protected Health Information (PHI) within encrypted on-premises database stores. Because no PHI is transmitted to external LLM vendors, enterprises are not required to execute Business Associate Agreements (BAAs) with third-party AI vendors.
\end{itemize}

\subsubsection{SOC 2 Type II and ISO/IEC 27001 Certification}
AMORA supports immutable append-only query audit logging, ensuring that every tool invocation, generated SQL query, and evidence score is permanently recorded with microsecond timestamps for internal security audits.

\subsection{Economic Cost Analysis and Break-Even Model}
To evaluate the commercial viability of AMORA, we construct an analytical Total Cost of Ownership (TCO) model comparing on-premises deployment against commercial cloud API endpoints.

\subsubsection{Cloud API Cost Formulation}
Let $Q$ denote the monthly query volume, let $\bar{T}_{\text{in}}$ denote the average prompt input tokens per query, and let $\bar{T}_{\text{out}}$ denote the average generated response tokens. Let $P_{\text{in}}$ and $P_{\text{out}}$ represent the commercial pricing per million tokens ($10^6$). The monthly operating expenditure (OpEx) of a cloud-based RAG deployment is given by:
\begin{equation}
    \text{OpEx}_{\text{Cloud}}(Q) = Q \times \left( \frac{\bar{T}_{\text{in}} \cdot P_{\text{in}}}{10^6} + \frac{\bar{T}_{\text{out}} \cdot P_{\text{out}}}{10^6} \right)
\end{equation}
For commercial frontier-tier APIs (e.g., standard enterprise rates of $\$3.00$ per $1\text{M}$ input tokens and $\$15.00$ per $1\text{M}$ output tokens), and adopting the empirical baseline prompt overhead of $\bar{T}_{\text{in}} = 1,073$ tokens (from Similarity-Only routing) with $\bar{T}_{\text{out}} = 150$ tokens:
\begin{equation}
    \begin{aligned}
        \text{Cost}_{\text{query, cloud}} &= \left( \frac{1{,}073 \times \$3.00}{10^6} \right) + \left( \frac{150 \times \$15.00}{10^6} \right) \\
        &= \$0.003219 + \$0.002250 = \$0.005469 \text{ per query}
    \end{aligned}
\end{equation}
For an enterprise processing $Q = 100,000$ queries per month:
\begin{equation}
    \text{OpEx}_{\text{Cloud}}(100{,}000) = 100{,}000 \times \$0.005469 = \mathbf{\$546.90 \text{ / month}}
\end{equation}
Over a 24-month operational horizon, cumulative cloud API expenditures total **\$13,125.60**.

\subsubsection{AMORA On-Premises Cost Formulation}
AMORA operates on a dedicated workstation GPU. Its economic cost structure comprises one-time capital expenditure (CapEx) amortized over its operational lifespan plus monthly electricity OpEx:
\begin{equation}
    \text{TCO}_{\text{AMORA}} = \frac{\text{CapEx}_{\text{Hardware}}}{T_{\text{lifespan}}} + \text{OpEx}_{\text{Electricity}} + \text{Cost}_{\text{Token}}
\end{equation}
\begin{itemize}
    \item \textbf{Hardware CapEx:} A dedicated NVIDIA GeForce RTX 3060 12GB workstation GPU entails a one-time acquisition cost of $\text{CapEx}_{\text{GPU}} = \$350.00$. Amortized over a standard $T_{\text{lifespan}} = 24$ months depreciation schedule:
    \begin{equation}
        \text{Amortization}_{\text{monthly}} = \frac{\$350.00}{24} = \mathbf{\$14.58 \text{ / month}}
    \end{equation}
    \item \textbf{Electricity Consumption OpEx:} The physical testbed was measured at an average active power draw of $P = 170\,\text{W}$ under continuous load. Assuming an enterprise operational profile of 8 active business hours per day, 22 business days per month:
    \begin{equation}
        \text{Energy}_{\text{monthly}} = 170\,\text{W} \times 8\,\text{hrs/day} \times 22\,\text{days} = 29.92\,\text{kWh/month}
    \end{equation}
    At commercial electricity tariffs of $\$0.15$ per $\text{kWh}$:
    \begin{equation}
        \text{OpEx}_{\text{Electricity}} = 29.92\,\text{kWh} \times \$0.15 = \mathbf{\$4.49 \text{ / month}}
    \end{equation}
    \item \textbf{Token Cost:} Because model weights are hosted locally via Ollama, marginal token API cost is identically zero:
    \begin{equation}
        \text{Cost}_{\text{Token}} = \mathbf{\$0.00}
    \end{equation}
\end{itemize}

Summing these components yields the total monthly operational cost of AMORA:
\begin{equation}
    \text{TCO}_{\text{AMORA}}(100{,}000) = \$14.58 + \$4.49 + \$0.00 = \mathbf{\$19.07 \text{ / month}}
\end{equation}

\subsubsection{Net Economic Savings and Payback Period}
Comparing AMORA against commercial cloud APIs:
\begin{equation}
    \text{Monthly Operational Savings} = \$546.90 - \$19.07 = \mathbf{\$527.83 \text{ / month}}
\end{equation}
\begin{equation}
    \text{Percentage Cost Reduction} = \frac{\$546.90 - \$19.07}{\$546.90} \times 100\% = \mathbf{96.51\%}
\end{equation}
The financial payback period required to fully amortize the initial GPU hardware investment is:
\begin{equation}
    \text{Payback Period} = \frac{\text{CapEx}_{\text{GPU}}}{\text{Monthly Net Savings}} = \frac{\$350.00}{\$527.83} \approx 0.66 \text{ months} \quad (\mathbf{\approx 20 \text{ calendar days}})
\end{equation}
By delivering a **$>96.5\%$ operational cost reduction** and achieving complete hardware payback in under three weeks, AMORA demonstrates overwhelming economic superiority for enterprise deployments.
